% page 400 of the facsimile (image p-399 of the scan)
% block 147.3 x 237.7 mm ; left margin 8.0 mm
\pgc{-12.700mm}{0.000mm}{
\l{\cel{2}392}
\l{}
\l{\sou{nektario}. (\sou{bot.}) Organaro glanda, situita en la floro e}
\l{\cel{3}qua kontenas suko mielatra. -}
\l{}
\l{\sou{nematoblasto}.\cel{2}(\sou{biol.}) Celulo specala, qua karakterizas}
\l{\cel{3}la celenterei, e qua kontenas organo urtikaganta. -}
\l{}
\l{\sou{nenio}. (\sou{epoki antiqua)}. Kanto funerala - da la ploristini,}
\l{\cel{3}lor la sepulto di personegi. - DEFI.}
\l{}
\l{\sou{nenufaro}. (\sou{bot.)}\cel{3}Planto aquala, tipo di la familio \textquotedbl{}nim-}
\l{\cel{3}feacei\textquotedbl{}. - L.\cel{1}\sou{Nuphar, Nymphae}a.- EFIS.}
\l{}
\l{\sou{neolamarckismo}.\cel{1}\sou{(biol.)}\cel{2}Teorio transformista qua admisas,}
\l{\cel{3}por explikar la origino di la speci, la influo prepon-}
\l{\cel{3}deranta, en la efekti da la uzado, di la efiki mekanis-}
\l{\cel{3}mala. - DEFIRS.}
\l{}
\l{\sou{neologismo.}\cel{2}Uzo di vorti nove-kreita, o di vorti anciena}
\l{\cel{3}ma uzata nova-senco. - DEFIS.}
\l{}
\l{\sou{neono}. (\sou{kemio}). Elemento qua esas gasa ye la temperaturo}
\l{\cel{3}ordinara, ed existas en aero minima-quante. -\cel{1}\sou{Simbolo}}
\l{\cel{3}\sou{kemiala}\cel{1}: Ne. - DEFIRS.}
\l{}
\l{\sou{neovitalismo}. (\sou{biol.}) Teorio segun qua la faktori fiziko-}
\l{\cel{3}kemiala qui kondicionas la fenomeni vivala efikas per}
\l{\cel{3}mekanismo partikulara e qui diferas de olta qua inter-}
\l{\cel{3}venas en la mondo materiala ordinara. - DEFIRS.}
\l{}
\l{\sou{neperala.}- (\sou{matem}.)Dicesas pri la logaritmi qui invent-}
\l{\cel{3}esis da Neper, la matematikisto Skota. -}
\l{}
\l{\sou{nepoto}. La filio di onua filio. - EFISL.}
\l{}
\l{\sou{nepotismo}. I. (\sou{historio eklez.})\cel{2}Favoro quan juis,lor ula}
\l{\cel{3}papi, la nevi e parenti di ilci. - II. Trouzo, da ofic-}
\l{\cel{3}iero, di sua povo\cel{2}por promocar neyuste sua familiani}
\l{\cel{3}o protektati. - DEFIS.}
\l{}
\l{\sou{neptunala}. (\sou{geol.).}\cel{1}Dicesas : 1. pri la tereni qui debas}
\l{\cel{3}sua origino ad aquo ; 2. pri sistemo qua explikis la}
\l{\cel{3}formaco di omna roki per la efiko da la aqui.}
\l{}
\l{\sou{nereido}. (\sou{mitol.}) Nimfo di maro.\cel{2}DEFIRS.}
\l{}
\l{\sou{nervo.}\cel{1}(\sou{anat.}) Singla de la filamenti qui igas la parti}
\l{\cel{3}diversa di la korpo homala komunikar kun la cerebro}
\l{\cel{3}e la myelo, transmisante ad olca la senti ; ad olta,}
\l{\cel{3}la impulsi movala. - DEFIRS.}
\l{}
\l{\sou{nervaturo.}\cel{1}I. (historio naturala) . Singla de la fileti}
\l{\cel{3}salianta qui quaze branchifas en la limbo di la folio,}
\l{\cel{3}che ula planti.\cel{18}(duras sur pag.393).}
}
